\documentclass[11pt]{article}\usepackage{layout}\begin{document}
\begin{canvas}\Header{第二章 · 全球开源生态地图}{teal}\Footer{01}
\Box{15}{27}{180}{\fontsize{20}{27}\selectfont\color{ink}\sffamily \bfseries 2.5 全球开源资产地图}
\Box{15}{40}{180}{\fontsize{9.7}{13.5}\selectfont\color{muted}\sffamily 代码、模型与数据共同构成人工智能时代的开放资源基础：既要看全球积累，也要看资源怎样进入实际使用。}
\Box{15}{53}{180}{\includegraphics[width=180mm]{.build/assets/fig-2-13-asset-world.pdf}}
\Box{15}{199}{180}{\fontsize{7.6}{10.7}\selectfont\color{muted}\sffamily 图 2-13　三类开源资产的全球分布与增长。代码为 GitHub 公开仓库；模型和数据总量来自 Hugging Face（HF），地域分布来自报告累计样本。来源：GitHub Innovation Graph、HF Summer 2026、报告观察数据；样本日期待核，口径详见附录 1。}
\Box{15}{212}{87}{\fontsize{10.3}{14.2}\selectfont\color{teal}\sffamily \bfseries ① 共同底座，三种不同的资源}
\Box{15}{218}{87}{\fontsize{9.9}{15}\selectfont\color{ink}代码承载可修改的软件，模型提供可调用、可改造的智能能力，数据支持训练、评测与适配。三类资源相互依赖；地图展示它们的供给位置，不把仓库数量相加为一个“资产总量”。}
\Box{108}{212}{87}{\fontsize{10.3}{14.2}\selectfont\color{teal}\sffamily \bfseries ② 同一国家，在不同资产上位置不同}
\Box{108}{218}{87}{\fontsize{9.9}{15}\selectfont\color{ink}美国在三类观察中均居前列；印度的公开代码仓库接近美国，中国在模型与数据样本中均居第二。代码积累的广度与 AI 资源供给的集中度并存，不能用单一平台替代整体生态判断。}
\Box{15}{246}{87}{\fontsize{10.3}{14.2}\selectfont\color{teal}\sffamily \bfseries ③ AI 时代的扩张是多层同步的}
\Box{15}{252}{87}{\fontsize{9.9}{15}\selectfont\color{ink}代码公开仓库存量一年增长 22.3\%；HF 模型、数据集在 2026 年 1—8 月分别增长约 21.8\% 和 40.6\%。时间跨度各异，但共同说明开放供给正在从软件代码扩展到智能能力与训练材料。}
\Box{108}{246}{87}{\fontsize{10.3}{14.2}\selectfont\color{teal}\sffamily \bfseries ④ 开放数量须进一步转化为公共能力}
\Box{108}{252}{87}{\fontsize{9.9}{15}\selectfont\color{ink}样本中，美国的数据集数量占 47.8\%，下载占 82.6\%，表明“有资源”与“被使用”并不等同。对中国与全球南方，资源建设还要连接本地需求、维护服务与全球分发，才能形成持续能力。}
\end{canvas}
\begin{canvas}\Header{第二章 · 全球开源生态地图}{teal}\Footer{02}
\Box{15}{27}{180}{\fontsize{20}{27}\selectfont\color{ink}\sffamily \bfseries 2.5.1 代码资产：从广泛积累到加速生产}
\Box{15}{40}{180}{\fontsize{9.7}{13.5}\selectfont\color{muted}\sffamily 公开仓库存量刻画积累，活跃仓库反映参与运转的基础，代码提交与公开贡献则揭示当期生产和协作的节奏。}
\Box{15}{53}{180}{\includegraphics[width=180mm]{.build/assets/fig-2-14-code-world.pdf}}
\Box{15}{199}{180}{\fontsize{7.6}{10.7}\selectfont\color{muted}\sffamily 图 2-14　公开代码仓库地域分布、增长与生产节奏。来源：GitHub Innovation Graph（2026 Q2）；OpenDigger 原报告平台活跃仓库序列；GitHub Octoverse 2026（2025.10—2026.09）。公开贡献与全部提交分列，不能相加。}
\Box{15}{212}{87}{\fontsize{10.3}{14.2}\selectfont\color{teal}\sffamily \bfseries ① 代码积累已形成广泛的地域基础}
\Box{15}{218}{87}{\fontsize{9.9}{15}\selectfont\color{ink}已披露地域的公开仓库合计 4.54 亿，约为 2020 年同期的 3.45 倍。报告年度 GitHub、Gitee 活跃仓库分别为 633 万、439 万：全球代码资源在积累，也通过不同平台持续运转。}
\Box{108}{212}{87}{\fontsize{10.3}{14.2}\selectfont\color{teal}\sffamily \bfseries ② 全球南方拥有可见的供给基础}
\Box{108}{218}{87}{\fontsize{9.9}{15}\selectfont\color{ink}印度、巴西、印度尼西亚分别有约 7,474 万、2,127 万和 1,138 万公开仓库，孟加拉国、尼日利亚也进入前二十。广泛的代码参与为本地应用与 AI 适配提供入口，其质量和持续维护仍需另行观察。}
\Box{15}{246}{87}{\fontsize{10.3}{14.2}\selectfont\color{teal}\sffamily \bfseries ③ 增长路径分化，不能只比较位次}
\Box{15}{252}{87}{\fontsize{9.9}{15}\selectfont\color{ink}美国与印度仓库存量接近，最近一年分别增长 17.0\% 和 35.0\%；巴西增长 24.7\%。中国合并内地及港澳台后约 3,388 万。地域数量受平台使用与归属方式影响，不足以单独判定技术领先。}
\Box{108}{246}{87}{\fontsize{10.3}{14.2}\selectfont\color{teal}\sffamily \bfseries ④ 生产加速，更需识别有效沉淀}
\Box{108}{252}{87}{\fontsize{9.9}{15}\selectfont\color{ink}2026 年度 GitHub 提交量达 343 亿次，公开贡献达 25 亿次。AI 正改变开发节奏，但一次提交不等于一项新资产；对高产出项目，还应联看问题解决、版本维护、测试与下游采用。}
\end{canvas}
\begin{canvas}\Header{第二章 · 全球开源生态地图}{teal}\Footer{03}
\Box{15}{27}{180}{\fontsize{20}{27}\selectfont\color{ink}\sffamily \bfseries 2.5.2 模型资产：供给扩张与全球再分发}
\Box{15}{40}{180}{\fontsize{9.7}{13.5}\selectfont\color{muted}\sffamily 模型仓库既包括原创模型，也包括微调、量化与格式转换等再创造；数量增长与下载流向共同揭示开放影响的形成。}
\Box{15}{53}{180}{\includegraphics[width=180mm]{.build/assets/fig-2-15-model-world.pdf}}
\Box{15}{199}{180}{\fontsize{7.6}{10.7}\selectfont\color{muted}\sffamily 图 2-15　模型供给地图、增长与下载结构。地域数量：报告发布者样本；增长：HF / AI World 日度序列及 HF Summer 2026；下载：Longpre 等《The Economies of Open Intelligence》（至 2025.08，滚动窗口过滤）。}
\Box{15}{212}{87}{\fontsize{10.3}{14.2}\selectfont\color{teal}\sffamily \bfseries ① 中美之外，欧洲与亚洲也有重要节点}
\Box{15}{218}{87}{\fontsize{9.9}{15}\selectfont\color{ink}样本中美国发布 45,801 个模型，中国发布 6,017 个；法国、德国、芬兰及日韩、新加坡各有供给。模型并非只由少数大厂生产，专业机构与社区共同连接基础模型、工具和应用。}
\Box{108}{212}{87}{\fontsize{10.3}{14.2}\selectfont\color{teal}\sffamily \bfseries ② 中国模型的使用位置高于数量份额}
\Box{108}{218}{87}{\fontsize{9.9}{15}\selectfont\color{ink}中国占发布者样本模型数量的 8.3\%，下载占 12.6\%；美国分别为 63.4\% 与 73.9\%。供给数量与使用集中度并不重合：中国的模型资源正获得超出其仓库数量份额的采用，仍需继续观察具体用途。}
\Box{15}{246}{87}{\fontsize{10.3}{14.2}\selectfont\color{teal}\sffamily \bfseries ③ 把模型变得可用，是重要的资产增量}
\Box{15}{252}{87}{\fontsize{9.9}{15}\selectfont\color{ink}HF 的 GGUF 格式仓库在 2026 年前七个月增长 464\%。这种格式便于模型在本地设备运行；微调让模型适配任务，量化通过压缩数值精度降低运行成本。开放的价值由原始权重延伸到这些再创造。}
\Box{108}{246}{87}{\fontsize{10.3}{14.2}\selectfont\color{teal}\sffamily \bfseries ④ 中国贡献在全球接续中得到放大}
\Box{108}{252}{87}{\fontsize{9.9}{15}\selectfont\color{ink}在 2024.08—2025.08 新发布模型的下载中，中国占 17.1\%，美国占 15.8\%；个人及跨地域组织合计 57.6\%。这是特定模型群体的下载结构，提示分发与再创造的重要性，不代表整体市场份额。}
\end{canvas}
\begin{canvas}\Header{第二章 · 全球开源生态地图}{teal}\Footer{04}
\Box{15}{27}{180}{\fontsize{20}{27}\selectfont\color{ink}\sffamily \bfseries 2.5.3 数据资产：从积累走向可持续使用}
\Box{15}{40}{180}{\fontsize{9.7}{13.5}\selectfont\color{muted}\sffamily 数据集提供训练、评测与本地化的材料。AI 将数据需求从文本扩展到动作、轨迹和环境交互，数据的可用性也因此更加重要。}
\Box{15}{53}{180}{\includegraphics[width=180mm]{.build/assets/fig-2-16-data-world.pdf}}
\Box{15}{199}{180}{\fontsize{7.6}{10.7}\selectfont\color{muted}\sffamily 图 2-16　数据集地域分布、增长与使用差异。地域数量和下载来自报告发布者样本；平台增长来自 HF Summer 2026；机器人数据集数量来自 HF Spring 2026。地图中未列明的地域不表示没有数据资产。}
\Box{15}{212}{87}{\fontsize{10.3}{14.2}\selectfont\color{teal}\sffamily \bfseries ① 数据供给的地域结构与模型不同}
\Box{15}{218}{87}{\fontsize{9.9}{15}\selectfont\color{ink}样本中，美国与中国分别有 5,871 和 1,338 个数据集；德国、法国之后，爱尔兰、马来西亚等也进入前列。较小经济体能够通过专业领域和语言资源形成贡献，不必复制大规模模型训练路线。}
\Box{108}{212}{87}{\fontsize{10.3}{14.2}\selectfont\color{teal}\sffamily \bfseries ② 中国有数量基础，还须增强使用连接}
\Box{108}{218}{87}{\fontsize{9.9}{15}\selectfont\color{ink}中国占样本数据集数量的 10.9\%，高于其模型数量份额 8.3\%，但数据集下载仅占 4.7\%。这一差异提示，应进一步检查语言覆盖、文档、格式、更新与分发条件，不能仅据此推断数据质量。}
\Box{15}{246}{87}{\fontsize{10.3}{14.2}\selectfont\color{teal}\sffamily \bfseries ③ 新增需求正在进入物理世界}
\Box{15}{252}{87}{\fontsize{9.9}{15}\selectfont\color{ink}HF 数据集总量在 2026 年 1—8 月从 71.1 万增至 100 万。机器人数据集从 2024 年的 1,145 个增至 2025 年的 26,991 个，约为前者的 23.6 倍；动作与交互记录正在成为训练智能系统的重要开放资源。}
\Box{108}{246}{87}{\fontsize{10.3}{14.2}\selectfont\color{teal}\sffamily \bfseries ④ 数据普惠依赖持续服务与本地参与}
\Box{108}{252}{87}{\fontsize{9.9}{15}\selectfont\color{ink}样本下载集中于美国，并不意味着需求只在那里。全球南方可围绕本地语言、农业、健康和公共服务建立数据合作；基金会等公共组织可支持明确使用条件、质量评测和维护，使资源持续可用。}
\end{canvas}
\begin{canvas}\Header{附录 1 · 2.5 口径与来源}{teal}\Footer{A1}
\Box{15}{28}{180}{\fontsize{20}{27}\selectfont\color{ink}\sffamily \bfseries 附录 1｜资产口径与证据来源}
\Box{15}{45}{180}{\fontsize{11}{16}\selectfont\color{teal}\sffamily \bfseries 1  本节中的“资产”与地域含义}
\Box{15}{54}{180}{\fontsize{9.3}{14.4}\selectfont\color{ink}本节以可见仓库观察资源供给：代码仓库是保存代码及协作记录的空间，模型仓库存放模型权重和配置，数据集仓库存放训练、评测等材料。公开可访问不自动等于符合开源许可；本节不把数量解释为技术质量、资产估值或开源合规数量。模型／数据归属发布单位，下载不代表使用者所在地。}
\Box{15}{84}{180}{\fontsize{11}{16}\selectfont\color{teal}\sffamily \bfseries 2  图中数量怎样取得与计算}
\Box{15}{93}{180}{\fontsize{9.3}{14.4}\selectfont\color{ink}代码地图：GitHub Innovation Graph 季末公开仓库存量，按具备 triage（事务分拣）及以上权限成员的地域众数归属；包括不活跃仓库。排除欧盟汇总行以免重复，中国合并内地、香港、澳门和台湾。4.54 亿为已披露地域合计；历史曲线均取第二季度。活跃仓库沿用报告原表，主期为 2025.07—2026.06；平台间不相加，也不与公开存量计算“活跃率”。}
\Box{15}{129}{180}{\fontsize{9.3}{14.4}\selectfont\color{ink}HF 地域样本：826 家发布单位、63 个国家与地区；原表分别列出 14 个地域，未列明不补零，快照日期与筛选细节待核。平台增长使用 HF 夏季材料的 1 月与 8 月端点，含期后观察；模型日度曲线截至 2025-12-08。下载结构比较“全历史模型”与 2024.08—2025.08 新发布模型，沿用原研究滚动窗口过滤，不能视为同口径年度市场份额变化。}
\Box{15}{164}{180}{\fontsize{11}{16}\selectfont\color{teal}\sffamily \bfseries 3  可追溯来源 · 2026-11-10 检索；点击条目访问原文}
\Box{15}{174}{180}{\sffamily\fontsize{8.9}{12}\selectfont S1\quad\href{https://github.com/github/innovationgraph/blob/main/data/repositories.csv}{GitHub Innovation Graph · repositories.csv / Datasheet}}
\Box{15}{181}{180}{\sffamily\fontsize{8.9}{12}\selectfont S2\quad\href{https://github.blog/news-insights/octoverse/octoverse-commits-pushes-reach-new-highs-in-2026-as-developers-build-with-agents/}{GitHub Octoverse 2026 · 提交、推送与公开贡献}}
\Box{15}{188}{180}{\sffamily\fontsize{8.9}{12}\selectfont S3\quad\href{https://github.com/will-ww/2026-GOSD-Report/tree/497478a19e92c676c225893a395738e747a71c8b/reports/round2-v2.0/source/chapters/02/modules/02/data/25}{GOSD 归档原表 · 活跃仓库、模型及数据地域样本}}
\Box{15}{195}{180}{\sffamily\fontsize{8.9}{12}\selectfont S4\quad\href{https://huggingface.co/blog/state-of-open-models-summer-2026}{Hugging Face · State of Open Models: Summer 2026 Observations}}
\Box{15}{202}{180}{\sffamily\fontsize{8.9}{12}\selectfont S5\quad\href{https://huggingface.co/blog/huggingface/state-of-os-hf-spring-2026}{Hugging Face · State of Open Source on Hugging Face: Spring 2026}}
\Box{15}{209}{180}{\sffamily\fontsize{8.9}{12}\selectfont S6\quad\href{https://aiworld.eu/story/hugging-faces-two-million-models-and-counting}{AI World / F. Ríos · Hugging Face’s Two Million Models and Counting}}
\Box{15}{216}{180}{\sffamily\fontsize{8.9}{12}\selectfont S7\quad\href{https://www.dataprovenance.org/economies-of-open-intelligence.pdf}{Longpre et al. · The Economies of Open Intelligence · Figure 1}}
\Box{15}{229}{180}{\fontsize{11}{16}\selectfont\color{teal}\sffamily \bfseries 4  v0.3 修订与本轮评审}
\Box{15}{238}{180}{\fontsize{9.3}{14.4}\selectfont\color{ink}本轮将原三页改为四页地图组合：三类资产总览、代码、模型、数据，共 16 个分析点。新增公开仓库地域底表、历史增长曲线、模型日度增长、下载群体比较；保留正文中的通俗解释。图号暂为 2-13—2-16，合稿时顺延下节。}
\Box{15}{263}{180}{\fontsize{9.3}{14.4}\selectfont\color{ink}待核查：HF 地域样本快照日期及筛选细节、活跃仓库的底层事件规则。评审重点：总览的三资产整合是否清晰；四页的证据与洞察是否足以支撑执行摘要及第六章。}
\end{canvas}\end{document}